\section{Introduction}

Quantitative real-time PCR (qPCR) is among the most widely used techniques
in molecular biology for quantifying nucleic acids and measuring gene
expression \cite{bustin2009miqe}. A sound qPCR analysis involves several
interdependent steps: estimating primer amplification efficiency, choosing
an appropriate quantification strategy, performing statistical testing,
and producing publication-quality figures \cite{bustin2010miqe}. When the
amplification efficiencies of target and reference genes are comparable,
the $2^{-\Delta\Delta C_T}$ method \cite{livak2001analysis} is convenient;
when they differ, an efficiency-corrected or standard-curve-based approach
is required instead \cite{pfaffl2001new}. These analytical choices, combined
with the need for rigorous statistics and clear visualization, make qPCR
data analysis demanding for many laboratory researchers.

A range of tools address qPCR analysis \cite{pabinger2014survey}, yet each
carries trade-offs. R packages---including our own \textit{qPCRtools}
\cite{li2022qpcrtools} and \textit{pcr} \cite{ahmed2018pcr}---implement
sound algorithms but demand R programming skill and a configured software
environment. Commercial packages such as QuantStudio and CFX Maestro offer
graphical interfaces but are costly, closed-source, and platform-restricted.
Free web-based tools are frequently discontinued and require uploading
experimental data to remote servers, raising reproducibility and privacy
concerns. The gap, therefore, is for a tool that combines
\emph{validated algorithmic equivalence} to a trusted implementation with a
\emph{native, cross-platform, zero-dependency} delivery.

We introduce \textit{qPCRtools}, a C++/Qt6 desktop application that
reimplements the analytical core of our R package behind a point-and-click
interface distributed as a single executable. This Application Note makes
four contributions: (i) a C++/Qt architecture in which a tested core
algorithm library is decoupled from a web-based user interface;
(ii) a validation demonstrating numerical equivalence to the R package
across all shared quantification methods; (iii) statistical and functional
extensions beyond the original R version; and (iv) a fully reproducible
validation package.
